% ============================================================
% MÓDULO: EL BANCO Y SUS PRODUCTOS
% ============================================================
\clearpage
\section*{MÓDULO: EL BANCO Y SUS PRODUCTOS}
\subsection*{Cómo manejarse como cliente}

\begin{recuadro}
\textbf{Objetivo del módulo:} conocer los productos bancarios más comunes, cómo se atiende a
un cliente dentro de un banco y cómo usar las cuentas, las tarjetas y los canales digitales de
forma segura y sin pagar intereses o comisiones innecesarias.
\end{recuadro}

% ------------------------------------------------------------
\section*{1. EL SISTEMA FINANCIERO}

\begin{itemize}
  \item \textbf{Bancos:} públicos (Banco Nación, Banco Provincia, etc.) y privados. Reciben
        depósitos, dan préstamos y prestan servicios de pago.
  \item \textbf{Banco Central (BCRA):} regula y controla a los bancos y protege a los
        usuarios de servicios financieros.
  \item \textbf{Billeteras virtuales y fintech:} aplicaciones que ofrecen cuentas virtuales
        (CVU), pagos con QR y transferencias. No son bancos, pero también están reguladas.
\end{itemize}

\paragraph*{Identificar una cuenta.}
\begin{itemize}
  \item \textbf{CBU} (Clave Bancaria Uniforme): 22 números que identifican una cuenta
        \textbf{bancaria}.
  \item \textbf{CVU} (Clave Virtual Uniforme): 22 números de una cuenta en una
        \textbf{billetera virtual}.
  \item \textbf{Alias:} palabras fáciles de recordar asociadas al CBU o CVU
        (\texttt{taller.avionica.quilmes}).
\end{itemize}
Las transferencias entre cuentas (bancarias o virtuales) son \textbf{inmediatas}, de día y de
noche.

% ------------------------------------------------------------
\section*{2. DENTRO DEL BANCO: QUIÉN ES QUIÉN}

\begin{center}
\small
\begin{tabular}{|p{0.25\linewidth}|p{0.65\linewidth}|}
\hline
\textbf{Sector o persona} & \textbf{Qué hace} \\ \hline
Atención al cliente / plataforma comercial & Abre cuentas, entrega tarjetas, resuelve
consultas y reclamos, ofrece productos. \\ \hline
\textbf{Oficial o ejecutivo de cuenta} & Es la persona asignada a un cliente (empresa,
comercio, profesional o cliente de mayores ingresos). Conoce su situación, le ofrece
productos, gestiona sus préstamos y le resuelve problemas. Es el contacto directo del cliente
con el banco. \\ \hline
\textbf{Caja} (ventanilla) & El \textbf{cajero} es el empleado que recibe depósitos en
efectivo, paga extracciones por ventanilla, cobra servicios y cambia moneda. Se atiende con
turno o número. \\ \hline
\textbf{Cajero automático} (ATM) & Máquina para extraer y depositar efectivo, transferir,
consultar saldos y pagar servicios con la tarjeta de débito, las 24 horas. \\ \hline
\textbf{Terminal de autoservicio} & Equipo para depositar cheques o efectivo, imprimir
resúmenes y hacer operaciones sin pasar por caja. \\ \hline
\end{tabular}
\end{center}

\paragraph*{Consejos para ir al banco.}
\begin{itemize}
  \item Sacar \textbf{turno} por la web o la aplicación (muchas sucursales lo exigen).
  \item Llevar el \textbf{DNI} y, si corresponde, la documentación pedida (recibo de sueldo,
        constancia de CUIT o de monotributo).
  \item Leer lo que se firma: el banco debe entregar copia del contrato con todas las
        condiciones y comisiones.
  \item Un banco \textbf{no puede obligar} a contratar otro producto (tarjeta de crédito,
        seguro, préstamo) como condición para abrir una caja de ahorro.
\end{itemize}

% ------------------------------------------------------------
\section*{3. LAS CUENTAS}

\begin{center}
\small
\begin{tabular}{|p{0.20\linewidth}|p{0.70\linewidth}|}
\hline
\textbf{Cuenta} & \textbf{Características} \\ \hline
\textbf{Caja de ahorro} & Para guardar dinero y operar. De apertura gratuita y sin costo de
mantenimiento, con tarjeta de débito. Puede ser en pesos o en dólares. \\ \hline
\textbf{Cuenta sueldo} & La abre el empleador para depositar el sueldo. Es gratuita para el
trabajador, incluida la tarjeta de débito. El trabajador puede elegir en qué banco cobrar.
Desde 2026 existe también en dólares. \\ \hline
\textbf{Cuenta corriente} & Para empresas y comercios. Permite emitir \textbf{cheques} y
puede tener \textbf{descubierto} (gastar más de lo que hay, pagando intereses). Tiene
comisiones. \\ \hline
\textbf{Cuenta en billetera virtual} & Se abre desde el celular, tiene CVU, permite cobrar,
pagar con QR y en muchos casos genera rendimientos diarios. \\ \hline
\end{tabular}
\end{center}

% ------------------------------------------------------------
\section*{4. TARJETA DE DÉBITO Y TARJETA DE CRÉDITO}

\begin{center}
\small
\begin{tabular}{|p{0.20\linewidth}|p{0.34\linewidth}|p{0.34\linewidth}|}
\hline
 & \textbf{Débito} & \textbf{Crédito} \\ \hline
¿De dónde sale el dinero? & De \textbf{tu cuenta}, en el momento. & Te lo \textbf{presta} el
banco; lo pagás después. \\ \hline
¿Límite? & El saldo de la cuenta. & Un límite de compra asignado según tus ingresos. \\ \hline
¿Cuotas? & No. & Sí (con o sin interés, según la promoción). \\ \hline
¿Intereses? & No. & Sí, si no pagás el total del resumen. \\ \hline
¿Costo? & Gratuita con caja de ahorro o cuenta sueldo. & Puede tener comisión de renovación
anual y mantenimiento. \\ \hline
\end{tabular}
\end{center}

\subsubsection*{4.1. Cómo funciona la tarjeta de crédito}
\begin{itemize}
  \item \textbf{Fecha de cierre:} último día en que las compras entran en el resumen del mes.
        Lo que se compra después del cierre se paga en el resumen siguiente.
  \item \textbf{Fecha de vencimiento:} día límite para pagar el resumen (unos días después del
        cierre).
  \item \textbf{Resumen:} detalla todas las compras, cuotas, impuestos y comisiones, el
        \textbf{saldo total} y el \textbf{pago mínimo}.
  \item \textbf{Opciones de pago:}
  \begin{itemize}
    \item \textbf{Pago total:} no se pagan intereses. \textbf{Es lo que hay que hacer siempre.}
    \item \textbf{Pago mínimo o parcial:} el resto se financia con \textbf{interés
          compensatorio}, que suele ser muy alto. La deuda crece mes a mes.
    \item \textbf{No pagar o pagar tarde:} se suman \textbf{intereses punitorios}, la deuda se
          informa a la central de deudores y se dificulta obtener créditos en el futuro.
  \end{itemize}
\end{itemize}

\begin{recuadro}
\textbf{Ejemplo:} un resumen de \$400.000 con pago mínimo de \$40.000. Si se paga solo el
mínimo, los \$360.000 restantes se financian a una tasa que puede superar el 100\% anual: en
un año la deuda puede casi duplicarse aunque no se compre nada más. \textbf{Regla de oro: la
tarjeta de crédito es un medio de pago, no un ingreso extra. Comprar solo lo que se puede pagar
completo al vencimiento.}
\end{recuadro}

\paragraph*{Seguridad con las tarjetas.}
Nunca compartir el PIN, el código de seguridad ni las claves; no prestar la tarjeta; avisar de
inmediato al banco si se pierde o la roban; revisar el resumen todos los meses y desconocer
compras que no se reconozcan.

% ------------------------------------------------------------
\section*{5. HOMEBANKING Y BANCA MÓVIL}

El \textbf{homebanking} (por computadora) y la \textbf{aplicación del banco} (en el celular)
permiten operar sin ir a la sucursal:
\begin{itemize}
  \item Consultar saldos, movimientos y resúmenes.
  \item Transferir y agendar destinatarios.
  \item Pagar servicios e impuestos, pagar la tarjeta de crédito, adherir débitos automáticos.
  \item Constituir plazos fijos, comprar dólares, pedir préstamos.
  \item Generar el \textbf{certificado de CBU} para darlo al empleador o a un cliente.
\end{itemize}

\subsubsection*{5.1. Seguridad digital: las estafas más comunes}
\begin{itemize}
  \item \textbf{Phishing:} correos, SMS o mensajes de WhatsApp que imitan al banco con un
        enlace falso para robar las claves.
  \item \textbf{Falso empleado del banco:} alguien llama diciendo que ``hay un problema con tu
        cuenta'' y pide códigos, el token o que vayas al cajero a ``validar'' algo.
  \item \textbf{Perfiles falsos en redes sociales} que ofrecen ayuda para reclamos.
\end{itemize}
\textbf{El banco nunca pide claves, PIN, token ni códigos por teléfono, mail o redes, y nunca
pide ir al cajero a hacer una operación.} Ante la duda, cortar y llamar al número que figura
en el dorso de la tarjeta.

% ------------------------------------------------------------
\section*{6. OTROS PRODUCTOS}

\begin{itemize}
  \item \textbf{Plazo fijo:} se deposita dinero por un plazo (30 días o más) a una tasa
        acordada. No se puede retirar antes.
  \item \textbf{Fondos comunes de inversión} de liquidez inmediata (``cuentas remuneradas''):
        rendimiento diario y disponibilidad inmediata.
  \item \textbf{Préstamos personales, prendarios e hipotecarios:} ver la unidad de préstamos
        (TNA, TEA y CFT).
  \item \textbf{Débito automático:} el banco paga una factura o la tarjeta de forma automática
        al vencimiento.
  \item \textbf{Seguros:} de vida, del hogar, de compras protegidas. Revisar si son
        obligatorios o se pueden dar de baja.
\end{itemize}

\paragraph*{Reclamos.}
Si el banco cobra algo que no corresponde, primero se reclama al banco (por escrito y con número
de reclamo). Si no responde en tiempo o la respuesta no conforma, se puede reclamar ante el
\textbf{BCRA} (Protección de los Usuarios de Servicios Financieros) o Defensa del Consumidor.

% ------------------------------------------------------------
\section*{7. ACTIVIDADES DEL MÓDULO}

\begin{enumerate}
  \item \textbf{Mi primera cuenta:} investigar qué requisitos piden dos bancos y una billetera
        virtual para abrir una cuenta. Comparar costos y servicios.
  \item \textbf{Leer un resumen:} con un resumen de tarjeta de crédito de ejemplo, identificar
        fecha de cierre, vencimiento, saldo total, pago mínimo, intereses e impuestos.
  \item \textbf{Simulación de deuda:} calcular en una planilla cómo evoluciona una deuda de
        tarjeta pagando solo el mínimo durante 12 meses.
  \item \textbf{Detectar el fraude:} analizar mensajes reales de phishing y explicar las señales
        de alerta.
  \item \textbf{Rol en el banco:} representar una consulta con el ejecutivo de cuenta (abrir una
        cuenta para un emprendimiento, pedir un préstamo para equipamiento).
\end{enumerate}
